% GENERATED FILE. Edit canonical lessons, then run npm run generate:books.
% canonical-source-hash: fnv1a64:5a2a0d050380eeb1
% canonical-lessons: MR-C297-ghusne, MR-C297-phodne, MR-C297-hiskavne, MR-C297-lutne, MR-C297-lubadne

\chapter{Thieves in the Night}
\label{ch:mr-thieves-in-the-night}
\hlchaptermodality{297}

\begin{chapteropening}
\textbf{By the end of this chapter:} \emph{I can tell how thieves broke in and what they took.}

Complete the last lesson of chapter 297: I can tell how thieves broke in and what they took.
\end{chapteropening}

\section[\texorpdfstring{\mr{घुसणे} (ghusṇe)}{घुसणे (ghusṇe)}]{\mr{घुसणे} (ghusṇe) — to barge in; to force one's way in}

\label{lesson:MR-C297-ghusne}



\begin{quote}
Before the new one: say the Marathi for to recover, then the Marathi for to be missed.
\end{quote}

\subsection*{You'll want to know: \mr{घुसणे}}
\textbf{\mr{घुसणे}} — \emph{ghusṇe} — \textquotedblleft{}to barge in; to force one's way in\textquotedblright{}.

\begin{quote}
\mr{रात्री} \mr{दोन} \mr{चोर} \mr{खिडकीतून} \mr{घरात} \mr{घुसले}.
\end{quote}

\begin{quote}
\emph{rātrī don cor khiḍkītūn gharāt ghusle.}
\end{quote}

\begin{quote}
At night two thieves forced their way into the house through the window.
\end{quote}

Past (he …): \mr{तो} \mr{घुसला}, \emph{to ghuslā}.

The place entered takes -\mr{त}: \mr{घरात} \mr{घुसणे}.

\begin{cousinweb}
\textbf{\mr{घुसणे}} (\emph{ghusṇe}) is inherited: CDIAL lists Marathi \emph{ghusṇe}, 'to enter forcibly', under \emph{*ghuss}, 'thrust in, pierce' (CDIAL 4492).
\end{cousinweb}

\begin{grammarlens}[title={What you've built}]
One more word for telling what happened.
\end{grammarlens}

\subsection*{Your turn}
\begin{itemize}
\raggedright
  \item \emph{Say it:} \emph{ghusṇe}
  \item \emph{Say it:} \emph{ghusṇe}, once more
  \item \emph{Say it:} \emph{rātrī don cor khiḍkītūn gharāt ghusle.} — At night two thieves forced their way into the house through the window.
  \item \emph{Recall:} say the Marathi for to recover, then the Marathi for to be missed, then say \emph{ghusṇe} again
\end{itemize}

\begin{tcolorbox}[breakable,colback=teal!4,colframe=teal!35!black,title={Before you move on}]
What does \mr{घुसणे} mean? (\textquotedblleft{}To barge in; to force one's way in\textquotedblright{}.) Say it once more.
\end{tcolorbox}
\section[\texorpdfstring{\mr{फोडणे} (phoḍṇe)}{फोडणे (phoḍṇe)}]{\mr{फोडणे} (phoḍṇe) — to break open; to smash}

\label{lesson:MR-C297-phodne}



\begin{quote}
Before the new one: say the Marathi for to be missed, then the Marathi for to barge in.
\end{quote}

\subsection*{You'll want to know: \mr{फोडणे}}
\textbf{\mr{फोडणे}} — \emph{phoḍṇe} — \textquotedblleft{}to break open; to smash\textquotedblright{}.

\begin{quote}
\mr{चोरांनी} \mr{कपाटाचं} \mr{कुलूप} \mr{फोडलं}.
\end{quote}

\begin{quote}
\emph{corānnī kapāṭācaṁ kulūp phoḍlaṁ.}
\end{quote}

\begin{quote}
The thieves broke the lock of the cupboard.
\end{quote}

Past, with -\mr{ने} (he …): \mr{त्याने} \mr{फोडलं}, \emph{tyāne phoḍlaṁ}.

Transitive: the thing broken agrees with the past: \mr{कुलूप} \mr{फोडलं}, \mr{खिडकी} \mr{फोडली}.

\begin{cousinweb}
\textbf{\mr{फोडणे}} (\emph{phoḍṇe}) is a tadbhava: CDIAL lists Marathi \emph{phoḍṇe}, 'to break', under Sanskrit \emph{sphoṭayati}, 'bursts, splits, breaks' (CDIAL 13857).
\end{cousinweb}

\begin{grammarlens}[title={What you've built}]
One more word for telling what happened.
\end{grammarlens}

\subsection*{Your turn}
\begin{itemize}
\raggedright
  \item \emph{Say it:} \emph{phoḍṇe}
  \item \emph{Say it:} \emph{phoḍṇe}, once more
  \item \emph{Say it:} \emph{corānnī kapāṭācaṁ kulūp phoḍlaṁ.} — The thieves broke the lock of the cupboard.
  \item \emph{Recall:} say the Marathi for to be missed, then the Marathi for to barge in, then say \emph{phoḍṇe} again
\end{itemize}

\begin{tcolorbox}[breakable,colback=teal!4,colframe=teal!35!black,title={Before you move on}]
What does \mr{फोडणे} mean? (\textquotedblleft{}To break open; to smash\textquotedblright{}.) Say it once more.
\end{tcolorbox}
\section[\texorpdfstring{\mr{हिसकावणे} (hiskāvṇe)}{हिसकावणे (hiskāvṇe)}]{\mr{हिसकावणे} (hiskāvṇe) — to snatch}

\label{lesson:MR-C297-hiskavne}



\begin{quote}
Before the new one: say the Marathi for to barge in, then the Marathi for to break open.
\end{quote}

\subsection*{You'll want to know: \mr{हिसकावणे}}
\textbf{\mr{हिसकावणे}} — \emph{hiskāvṇe} — \textquotedblleft{}to snatch\textquotedblright{}.

\begin{quote}
\mr{एका} \mr{माणसाने} \mr{तिच्या} \mr{हातातली} \mr{पिशवी} \mr{हिसकावली}.
\end{quote}

\begin{quote}
\emph{ekā māṇsāne ticyā hātātlī piśvī hiskāvlī.}
\end{quote}

\begin{quote}
A man snatched the bag from her hand.
\end{quote}

Past, with -\mr{ने} (he …): \mr{त्याने} \mr{हिसकावलं}, \emph{tyāne hiskāvlaṁ}.

The thing snatched is the object and agrees with the past: \mr{पिशवी} \mr{हिसकावली}.

\begin{cousinweb}
\textbf{\mr{हिसकावणे}} (\emph{hiskāvṇe}) is inherited: it belongs to the family CDIAL lists under \emph{*hicc} (CDIAL 14083), where Marathi \emph{hisakṇe} is 'to jerk'. To snatch is to jerk away.
\end{cousinweb}

\begin{grammarlens}[title={What you've built}]
One more word for telling what happened.
\end{grammarlens}

\subsection*{Your turn}
\begin{itemize}
\raggedright
  \item \emph{Say it:} \emph{hiskāvṇe}
  \item \emph{Say it:} \emph{hiskāvṇe}, once more
  \item \emph{Say it:} \emph{ekā māṇsāne ticyā hātātlī piśvī hiskāvlī.} — A man snatched the bag from her hand.
  \item \emph{Recall:} say the Marathi for to barge in, then the Marathi for to break open, then say \emph{hiskāvṇe} again
\end{itemize}

\begin{tcolorbox}[breakable,colback=teal!4,colframe=teal!35!black,title={Before you move on}]
What does \mr{हिसकावणे} mean? (\textquotedblleft{}To snatch\textquotedblright{}.) Say it once more.
\end{tcolorbox}
\section[\texorpdfstring{\mr{लुटणे} (luṭṇe)}{लुटणे (luṭṇe)}]{\mr{लुटणे} (luṭṇe) — to loot; to plunder}

\label{lesson:MR-C297-lutne}



\begin{quote}
Before the new one: say the Marathi for to break open, then the Marathi for to snatch.
\end{quote}

\subsection*{You'll want to know: \mr{लुटणे}}
\textbf{\mr{लुटणे}} — \emph{luṭṇe} — \textquotedblleft{}to loot; to plunder\textquotedblright{}.

\begin{quote}
\mr{रस्त्यात} \mr{चार} \mr{माणसांनी} \mr{प्रवाशांना} \mr{लुटलं}.
\end{quote}

\begin{quote}
\emph{rastyāt cār māṇsānnī pravāśānnā luṭlaṁ.}
\end{quote}

\begin{quote}
On the road four men robbed the travellers.
\end{quote}

Past, with -\mr{ने} (he …): \mr{त्याने} \mr{लुटलं}, \emph{tyāne luṭlaṁ}.

The victims take -\mr{ला}: \mr{प्रवाशांना} \mr{लुटणे}; goods are the object: \mr{दुकान} \mr{लुटणे}.

\begin{cousinweb}
\textbf{\mr{लुटणे}} (\emph{luṭṇe}) is inherited: CDIAL lists Marathi \emph{luṭṇe}, 'plunders', under \emph{*luṭṭati}, 'plunders' (CDIAL 11078).
\end{cousinweb}

\begin{grammarlens}[title={What you've built}]
One more word for telling what happened.
\end{grammarlens}

\subsection*{Your turn}
\begin{itemize}
\raggedright
  \item \emph{Say it:} \emph{luṭṇe}
  \item \emph{Say it:} \emph{luṭṇe}, once more
  \item \emph{Say it:} \emph{rastyāt cār māṇsānnī pravāśānnā luṭlaṁ.} — On the road four men robbed the travellers.
  \item \emph{Recall:} say the Marathi for to break open, then the Marathi for to snatch, then say \emph{luṭṇe} again
\end{itemize}

\begin{tcolorbox}[breakable,colback=teal!4,colframe=teal!35!black,title={Before you move on}]
What does \mr{लुटणे} mean? (\textquotedblleft{}To loot; to plunder\textquotedblright{}.) Say it once more.
\end{tcolorbox}
\section[\texorpdfstring{\mr{लुबाडणे} (lubāḍṇe)}{लुबाडणे (lubāḍṇe)}]{\mr{लुबाडणे} (lubāḍṇe) — to fleece; to swindle}

\label{lesson:MR-C297-lubadne}



\begin{quote}
Before the new one: say the Marathi for to snatch, then the Marathi for to loot.
\end{quote}

\subsection*{You'll want to know: \mr{लुबाडणे}}
\textbf{\mr{लुबाडणे}} — \emph{lubāḍṇe} — \textquotedblleft{}to fleece; to swindle\textquotedblright{}.

\begin{quote}
\mr{त्या} \mr{माणसाने} \mr{गोड} \mr{बोलून} \mr{आजीला} \mr{लुबाडलं}.
\end{quote}

\begin{quote}
\emph{tyā māṇsāne goḍ bolūn ājīlā lubāḍlaṁ.}
\end{quote}

\begin{quote}
That man swindled Grandma with sweet talk.
\end{quote}

Past, with -\mr{ने} (he …): \mr{त्याने} \mr{लुबाडलं}, \emph{tyāne lubāḍlaṁ}.

Like \mr{लुटणे}, but by trickery; the victim takes -\mr{ला}.

\begin{cousinweb}
\textbf{\mr{लुबाडणे}} (\emph{lubāḍṇe}) is a Marathi verb of fleecing by trickery; its origin is not settled.
\end{cousinweb}

\begin{grammarlens}[title={What you've built}]
One more word for telling what happened.
\end{grammarlens}

\subsection*{Your turn}
\begin{itemize}
\raggedright
  \item \emph{Say it:} \emph{lubāḍṇe}
  \item \emph{Say it:} \emph{lubāḍṇe}, once more
  \item \emph{Say it:} \emph{tyā māṇsāne goḍ bolūn ājīlā lubāḍlaṁ.} — That man swindled Grandma with sweet talk.
  \item \emph{Recall:} say the Marathi for to snatch, then the Marathi for to loot, then say \emph{lubāḍṇe} again
\end{itemize}

\begin{tcolorbox}[breakable,colback=teal!4,colframe=teal!35!black,title={Before you move on}]
What does \mr{लुबाडणे} mean? (\textquotedblleft{}To fleece; to swindle\textquotedblright{}.) Say it once more.
\end{tcolorbox}
